% ============================================================
%  body.tex - Experiment A: OpenWAM-alpha on LIBERO
%  Section-only content (no preamble). Imported by
%      reports/experiments/2026.10.05-2026.10.11-openwam-libero/openwam-libero.tex
%      report.tex
%  The importing document must load the `ego-report` style first.
% ============================================================

\section{Experiment A: \texorpdfstring{OpenWAM-$\alpha$}{OpenWAM-alpha} on LIBERO}
\label{sec:owa}

\subsection{Overview and Scope}
\label{sec:owa-overview}

This experiment reproduces the published LIBERO success rate of the
\emph{OpenWAM-$\alpha$} policy~\cite{openwam2026} on a single A100 GPU. No new
weights were trained and no per-seed tuning was performed: the authors'
released LIBERO checkpoint is evaluated under its native task protocol, and the
resulting score is compared against the officially reported number.

LIBERO contains 130 language-conditioned manipulation tasks in total. This
evaluation covers the four standard suites --- Spatial, Object, Goal and
LIBERO-10 --- that together contain 40 tasks. Each task is rolled out for 50
episodes, giving $40 \times 50 = 2000$ closed-loop episodes. The complete task
definitions, scene assets and fixed initial states are taken from a pinned
version of LIBERO. Expert demonstrations are not required for evaluation and
the authors' full training set was not downloaded.

\subsection{The LIBERO Benchmark}
\label{sec:owa-benchmark}

LIBERO is a language-conditioned manipulation benchmark for simulated robot
arms~\cite{libero2023}. Every task defines a tabletop scene, a natural-language
goal and a BDDL goal predicate. At each control step the policy observes the
current external and wrist RGB frames together with the robot's proprioceptive
state, predicts an action, and the MuJoCo simulator advances the scene to
produce the next observation. This is a closed loop, not video classification
and not a real-robot deployment.

The raw data therefore consists of scene assets, BDDL task descriptions,
language instructions and fixed initial states. RGB frames are rendered
dynamically from the geometry, joint, mass, friction, contact and camera
configuration of the simulator; the model is never handed a pre-recorded
future video. Object ground-truth positions and the success predicate are
maintained by the environment and are not extra policy inputs.

A few terms are used with a precise meaning throughout this report:

\begin{itemize}
  \item \textbf{task} --- the goal, for example \emph{``put the milk in the
        basket''};
  \item \textbf{trial / episode} --- one complete attempt starting from a fixed
        initial state;
  \item \textbf{control step} --- one execution of a 7-dimensional arm command;
  \item \textbf{query} --- one network planning call, which may emit several
        actions at once.
\end{itemize}

Trial IDs 0--49 index the 50 attempts of a task and select the initial state;
they are \emph{not} seeds. A seed is the origin of a random-number stream;
fixing it aids reproducibility but is not by itself sufficient --- initial
states, library versions, numerical precision and model stochasticity must be
fixed as well.

The final LIBERO action is a 3-dimensional translation command, a
3-dimensional rotation command and a 1-dimensional gripper command. It is not
a set of seven joint angles, and the numbers cannot be read directly as metres.
At a 20\,Hz control rate each step corresponds to roughly $0.05$\,s of
simulated time; internal physics sub-steps are not counted as policy steps.
A budget of 600 or 700 steps therefore corresponds to about 30 or 35 seconds of
simulated time, which is unrelated to wall-clock runtime. The budget is the
maximum number of control steps allowed --- not a monetary or GPU-capacity
limit. An episode stops as soon as the goal is satisfied; if the budget is
exhausted the episode is recorded as a failure. The initial stabilisation wait
is not charged against the policy budget.

Success is measured by the task-specific scoring rule, which inspects spatial
relations, containment relations and joint or switch states of the objects.
The policy sees the natural-language goal while the scorer reads the
machine-checkable goal predicate. Success rate is defined as successful
episodes divided by all episodes; with 50 trials per task and 500 per suite,
and with equal denominators across the four suites, the mean of the suite
scores equals the overall 2000-episode success rate. All failures are retained
and engineering errors are reported separately; incomplete runs are never
presented as full results.

Official summary scores are not a per-episode ground-truth archive. Agreement
with the official number does not imply sample-by-sample identity, and because
the two models reproduced in this report use different native budgets, seeds,
simulator versions and replanning rates, the difference between their scores
cannot be interpreted as a fair architectural advantage.

\subsection{Model and Inference Architecture}
\label{sec:owa-model}

The checkpoint used here is
\btt{OpenWAM-Alpha-Sim-LIBERO/checkpoint\_step\_10690.safetensors},
pinned to Hugging Face revision
\btt{110b07e4d5581c81520a77419cbb0e9371df231b}.
All nine files (24{,}835{,}230{,}170 bytes in total) were verified; the main
weight file is 24{,}813{,}767{,}464 bytes with SHA-256
\btt{34c9a6c9af67dda9aaffa2edc0efa29b28eeee1c26dea419f98065c80508c96a}.

The loading log confirms a Wan video DiT backbone, the WanVideoVAE38 decoder, a
Wan text encoder and a dedicated action expert, composed through the authors'
\emph{Dual-System Joint Self-Attention} with \texttt{bfloat16} mixed precision.
A Wan2.2-TI2V-5B video base~\cite{wan2025} and the action branch are combined by
joint self-attention, so that a video-generation backbone participates in
world/action modelling --- a design that differs from the vision-language
backbone plus action expert used by $\pi_{0.5}$ (Experiment~B). The dual-system
joint modelling and the unified action representation are the authors'
contributions, not contributions of this report.

Internally the model works in an 80-dimensional unified action representation.
A normaliser maps this representation to and from a native 10-dimensional
end-effector action, and the client converts the rotation part into LIBERO's
7-dimensional control, so the first seven entries cannot simply be sliced off.

Figure~\ref{fig:owa-dataflow} lists the runtime data flow, which is a
deployment diagram rather than a layer-by-layer network topology.

\begin{figure}[H]
  \centering
  \fbox{\parbox{0.92\linewidth}{\centering\small
    task language $\rightarrow$ author tokenizer / text encoder;\\
    external + wrist RGB $\rightarrow$ two-view canvas / video encoder;\\
    EEF 10-D state $\rightarrow$ min-max normalisation / 80-D mapping;\\
    (text, video, state) $\rightarrow$ joint denoising of video backbone and
    action expert;\\
    $\rightarrow$ native EEF 10-D action $\rightarrow$ rotation conversion
    $\rightarrow$ 7-D LIBERO control;\\
    $\rightarrow$ execute and re-observe (feedback loop).}}
  \caption{Runtime data flow of the OpenWAM-$\alpha$ closed loop.}
  \label{fig:owa-dataflow}
\end{figure}

The server logs report a $384 \times 320$ canvas with multi-view enabled; two
cameras are actually present and the third slot is empty. Denoising uses a
10-step schedule in synchronous mode with a deployment horizon of 10. The
authors additionally enable a DiT velocity cache with a threshold of 0.990, so
the ten denoising steps must not be described as ten independent full video-DiT
forward passes.

\emph{What is evaluated here is action success rate, not future-video quality.}
The full run stored no predicted future RGB and no future ground truth, so no
world-model video metric, counterfactual accuracy or causal effect is reported.
The fact that the architecture is a world--action model does not by itself
demonstrate such capabilities in this experiment.

\subsection{Closed-Loop Data Pipeline}
\label{sec:owa-pipeline}

A single episode proceeds as follows.

\begin{enumerate}
  \item Task definitions, language, scene assets and initial states are read
        from the pinned LIBERO repository. Judging the released weights does not
        require the full training demonstration set, and this run did not
        download it.
  \item One simulator instance is created per task and trials 0--49 are executed
        consecutively; after each reset the corresponding initial state is
        applied. The environment seed is 42.
  \item Thirty stabilisation steps are executed and are not charged against the
        policy budget (600 for Spatial/Object/Goal, 700 for LIBERO-10).
  \item External and wrist RGB images and the robot state are captured. The
        author-side end-effector state is a 3-D position, a continuous 6-D
        rotation and a 1-D gripper value. Object ground-truth fields exist in
        the simulator observation but were not fed to the policy.
  \item The client packages images, task and state and sends the PNG payload
        over a WebSocket to a separate model process. The teaching hook records
        the genuine PNG bytes; no camera layout or $\pi_{0.5}$-style state
        adaptation is substituted.
  \item The server applies min-max normalisation from the checkpoint's
        \texttt{normalization\_stats.npy}, runs the joint denoising, and hands
        the output to the adapter.
  \item The client converts the 6-D rotation of the native 10-D action into
        LIBERO rotation control to form a 7-D action and executes it at the
        author's caching cadence; typically a 10-step chunk is consumed before
        replanning. The 74-control-step teaching case corresponds to 8 planning
        calls --- control steps are not neural inference calls.
  \item Each action changes the scene and the next observation re-enters the
        loop. \texttt{env.done} is the success signal; episodes that time out are
        still counted in the denominator.
  \item Per-task \texttt{results.json} files and the client \texttt{RESULT} log
        store every trial, its success flag, policy steps and terminal reward;
        the run is then summarised into \texttt{summary.json} and
        \texttt{summary.csv}.
\end{enumerate}

Debug logging was disabled for the full run. Ninety log, configuration and
result files were retained, but per-query PNGs, action chunks and states were
not stored for the main table. A single teaching case additionally stores 74
genuine observation/action frames; it is a learning exhibit and does not
replace the 2000-episode result.

\subsection{Frozen Environment and Provenance}
\label{sec:owa-env}

\begin{table}[H]
  \centering
  \small
  \caption{Frozen configuration of the full OpenWAM-$\alpha$ run.}
  \label{tab:owa-env}
  \begin{tabular}{@{}p{0.30\linewidth}p{0.64\linewidth}@{}}
    \toprule
    Item & Value \\
    \midrule
    Author source & \btt{48bd67b89d489b14d03b8d92bc66e65d306df32e} \\
    LIBERO source & \btt{8f1084e3132a39270c3a13ebe37270a43ece2a01} \\
    Python (model/client) & 3.11.16 \\
    PyTorch (model) & 2.5.1+cu121 \\
    Transformers & 5.5.0 \\
    NumPy & 1.26.4 \\
    MuJoCo / robosuite & 3.3.2 / 1.4.0 \\
    dtype / attention & \texttt{bfloat16} / Torch SDPA \\
    \texttt{compile} & false \\
    TF32 & off, precision \emph{highest}, \btt{NVIDIA\_TF32\_OVERRIDE=0} \\
    Hardware & single A100, EGL, model and simulator in separate processes \\
    Model server replicas & 1, port 8920, concurrency unchanged mid-run \\
    Environment seed / trial & 42 / 0--49 per task \\
    Checkpoint saving & none (inference only, no new weights) \\
    \bottomrule
  \end{tabular}
\end{table}

The authors recommend Python 3.10 and Torch 2.7; this run reuses an existing
Python 3.11 / Torch 2.5.1 environment with SDPA and \texttt{compile} disabled.
These deviations are disclosed and no claim of bitwise equality, equal speed or
a fully reproduced author environment is made. System drivers, CUDA, the base
Conda environment and the global shell were left untouched.

The formal run ID is \btt{openwam-w2-libero-native-full-s42}. The
pre-registered source Git revision was
\btt{0c6e95571585a29914f62e7036768be57379482d}
and the queue file SHA-256 was
\btt{1078f11d00808fe962c6215e03f577499e219136dcb0d9cf6b1140bb495a25f}.
The executed launch command was:

\begin{verbatim}
server-python OpenWAM/benchmarks/libero/run_all_suites.py \
  --ckpt-dir <models>/OpenWAM-Alpha-Sim-LIBERO \
  --ckpt-name checkpoint_step_10690.safetensors \
  --server-python <server-python> --libero-python <client-python> \
  --libero-path <LIBERO-8f1084e3...> \
  --policy-config <OpenWAM>/benchmarks/libero/policy_config.yml \
  --gpus 0 --render-gpus 0 --replicas-per-gpu 1 \
  --compile-enabled false --client-max-attempts 1 \
  --output-dir <runs>/openwam-w2-libero-native-full-s42/output
\end{verbatim}

Re-evaluating the policy requires a new unique run ID, a fresh space/ETA
registration and a new configuration; this output must not be overwritten.

\subsection{Results}
\label{sec:owa-results}

The full 2000-episode run completed with 1985 successes and 15 failures, giving
an overall success rate of 99.25\,\%. Table~\ref{tab:owa-scores} breaks the
result down by suite and compares it with the authors' reported numbers.

\begin{table}[H]
  \centering
  \small
  \caption{Per-suite success rate: this reproduction versus the official
  OpenWAM-$\alpha$ scores.}
  \label{tab:owa-scores}
  \begin{tabular}{@{}lrrrr@{}}
    \toprule
    Suite & Success / total & This run & Official & Difference (pp) \\
    \midrule
    libero\_spatial & 497 / 500 & 99.40\,\% & 99.60\,\% & $-0.20$ \\
    libero\_object  & 498 / 500 & 99.60\,\% & 99.60\,\% & $+0.00$ \\
    libero\_goal    & 497 / 500 & 99.40\,\% & 99.80\,\% & $-0.40$ \\
    libero\_10      & 493 / 500 & 98.60\,\% & 98.20\,\% & $+0.40$ \\
    \midrule
    \textbf{Mean / total} & \textbf{1985 / 2000} & \textbf{99.25\,\%} &
    \textbf{99.30\,\%} & $\mathbf{-0.05}$ \\
    \bottomrule
  \end{tabular}
\end{table}

The official table publishes only summary scores and provides no per-episode
predictions, so sample-by-sample agreement with the authors cannot be claimed.
The value 99.25\,\% is reported to two decimals and is not described as exactly
matching 99.30\,\% merely because of rounding. A single fixed seed cannot
demonstrate multi-seed stability, statistical superiority or a new method.
Figure~\ref{fig:owa-official} shows the aggregate comparison and
Figure~\ref{fig:owa-40tasks} the per-task breakdown.

\begin{figure}[H]
  \centering
  \includegraphics[width=0.85\linewidth]{2026.10.05-2026.10.11-openwam-libero/official-comparison.png}
  \caption{Official OpenWAM-$\alpha$ scores versus this reproduction.}
  \label{fig:owa-official}
\end{figure}

\begin{figure}[H]
  \centering
  \includegraphics[width=0.95\linewidth]{2026.10.05-2026.10.11-openwam-libero/all-40-tasks.png}
  \caption{Per-task success rate over all 40 LIBERO tasks.}
  \label{fig:owa-40tasks}
\end{figure}

Table~\ref{tab:owa-40tasks} lists every task together with the number of
successful episodes and the failing trial IDs.

\begin{table}[H]
  \centering
  \footnotesize
  \caption{Per-task results for the full OpenWAM-$\alpha$ run.}
  \label{tab:owa-40tasks}
  \begin{tabular}{@{}l r p{7.0cm} r l@{}}
    \toprule
    Suite & ID & Task language & Succ./50 & Failing trials \\
    \midrule
    libero\_spatial & 0 & pick up the black bowl between the plate and the ramekin and place it on the plate & 50 & --- \\
    libero\_spatial & 1 & pick up the black bowl next to the ramekin and place it on the plate & 50 & --- \\
    libero\_spatial & 2 & pick up the black bowl from table center and place it on the plate & 50 & --- \\
    libero\_spatial & 3 & pick up the black bowl on the cookie box and place it on the plate & 50 & --- \\
    libero\_spatial & 4 & pick up the black bowl in the top drawer of the wooden cabinet and place it on the plate & 48 & 4, 27 \\
    libero\_spatial & 5 & pick up the black bowl on the ramekin and place it on the plate & 49 & 2 \\
    libero\_spatial & 6 & pick up the black bowl next to the cookie box and place it on the plate & 50 & --- \\
    libero\_spatial & 7 & pick up the black bowl on the stove and place it on the plate & 50 & --- \\
    libero\_spatial & 8 & pick up the black bowl next to the plate and place it on the plate & 50 & --- \\
    libero\_spatial & 9 & pick up the black bowl on the wooden cabinet and place it on the plate & 50 & --- \\
    libero\_object & 0 & pick up the alphabet soup and place it in the basket & 50 & --- \\
    libero\_object & 1 & pick up the cream cheese and place it in the basket & 50 & --- \\
    libero\_object & 2 & pick up the salad dressing and place it in the basket & 50 & --- \\
    libero\_object & 3 & pick up the bbq sauce and place it in the basket & 50 & --- \\
    libero\_object & 4 & pick up the ketchup and place it in the basket & 50 & --- \\
    libero\_object & 5 & pick up the tomato sauce and place it in the basket & 50 & --- \\
    libero\_object & 6 & pick up the butter and place it in the basket & 50 & --- \\
    libero\_object & 7 & pick up the milk and place it in the basket & 48 & 10, 13 \\
    libero\_object & 8 & pick up the chocolate pudding and place it in the basket & 50 & --- \\
    libero\_object & 9 & pick up the orange juice and place it in the basket & 50 & --- \\
    libero\_goal & 0 & open the middle drawer of the cabinet & 50 & --- \\
    libero\_goal & 1 & put the bowl on the stove & 50 & --- \\
    libero\_goal & 2 & put the wine bottle on top of the cabinet & 48 & 38, 47 \\
    libero\_goal & 3 & open the top drawer and put the bowl inside & 50 & --- \\
    libero\_goal & 4 & put the bowl on top of the cabinet & 50 & --- \\
    libero\_goal & 5 & push the plate to the front of the stove & 50 & --- \\
    libero\_goal & 6 & put the cream cheese in the bowl & 50 & --- \\
    libero\_goal & 7 & turn on the stove & 50 & --- \\
    libero\_goal & 8 & put the bowl on the plate & 49 & 25 \\
    libero\_goal & 9 & put the wine bottle on the rack & 50 & --- \\
    libero\_10 & 0 & put both the alphabet soup and the tomato sauce in the basket & 47 & 7, 14, 38 \\
    libero\_10 & 1 & put both the cream cheese box and the butter in the basket & 50 & --- \\
    libero\_10 & 2 & turn on the stove and put the moka pot on it & 50 & --- \\
    libero\_10 & 3 & put the black bowl in the bottom drawer of the cabinet and close it & 50 & --- \\
    libero\_10 & 4 & put the white mug on the left plate and put the yellow and white mug on the right plate & 50 & --- \\
    libero\_10 & 5 & pick up the book and place it in the back compartment of the caddy & 50 & --- \\
    libero\_10 & 6 & put the white mug on the plate and put the chocolate pudding to the right of the plate & 49 & 23 \\
    libero\_10 & 7 & put both the alphabet soup and the cream cheese box in the basket & 50 & --- \\
    libero\_10 & 8 & put both moka pots on the stove & 48 & 29, 38 \\
    libero\_10 & 9 & put the yellow and white mug in the microwave and close it & 49 & 25 \\
    \bottomrule
  \end{tabular}
\end{table}

Independent verification read all 40 \texttt{results.json} files and matched
each task against the 50 client \texttt{RESULT} entries on trial ID, success
flag and policy steps, additionally checking seed, stabilisation steps, budget,
end-effector mode, coverage, reward, summary and manifest. All 2000 unique
trials, 1985 successes and 15 failures agree. The 90 formal files
(10{,}271{,}072 bytes) were returned with per-file size/SHA checks and a server
acknowledgement fingerprint of
\btt{ca759e55b79bb6c7556399a685662773481eb07f91d908fd62c6d11d4cf91f21}.
Every complete episode retains only its task, trial, outcome, action count and
terminal reward --- not the full action sequence, frame-by-frame media or
environment state.

\subsection{Compute, Timing and Control Budget}
\label{sec:owa-compute}

The formal run started at 2026-10-05T18:40:58+08:00 and finished at
2026-10-06T01:51:27+08:00 with exit code 0. The total wall time was
25{,}829.198 seconds (7\,h\,10\,m\,29\,s), including model loading, compilation
and caching, simulation, serialisation, task switching and log writing, but
excluding model download, earlier failed attempts and document preparation.

Full-device sampling at 2-second intervals yielded 12{,}617 points, with a peak
memory of 25{,}676\,MiB ($\approx 25.07$\,GiB) and an average GPU utilisation of
29.40\,\%. A momentary high utilisation does not mean the GPU was saturated for
the whole run, and the sampled peak may miss short spikes between samples.

\begin{figure}[H]
  \centering
  \includegraphics[width=0.85\linewidth]{2026.10.05-2026.10.11-openwam-libero/gpu-resource-history.png}
  \caption{Measured GPU memory and utilisation over the full run.}
  \label{fig:owa-gpu}
\end{figure}

A single model replica used about 25\,GiB, which suggests that two replicas
would fit on the device, but whether that would reduce wall time still requires
a controlled measurement under a fixed protocol; a linear speed-up is not
promised. Any future increase in concurrency must record the model's random
streams and task assignment and must not silently change the replica count of
this frozen run.

\begin{table}[H]
  \centering
  \small
  \caption{Executed control steps per suite.}
  \label{tab:owa-steps}
  \begin{tabular}{@{}lrr@{}}
    \toprule
    Suite & Total control steps & Median per episode \\
    \midrule
    libero\_spatial & 51{,}855  & 100.5 \\
    libero\_object  & 66{,}294  & 129.0 \\
    libero\_goal    & 56{,}520  & 92.0  \\
    libero\_10      & 127{,}141 & 238.0 \\
    \midrule
    \textbf{Total}  & \textbf{301{,}810} & --- \\
    \bottomrule
  \end{tabular}
\end{table}

The 301{,}810 executed control steps are not equal to the number of model
queries; the full run has no per-query audit, so no exact call count or
pure-GPU throughput is invented.

\subsection{Failure Inventory}
\label{sec:owa-failures}

All 15 failures are listed in Table~\ref{tab:owa-failures}. Each exhausted its
budget (600 steps for Spatial/Object/Goal, 700 for LIBERO-10) and ended with a
terminal reward of 0.0.

\begin{table}[H]
  \centering
  \small
  \caption{All 15 failed episodes of the full run.}
  \label{tab:owa-failures}
  \begin{tabular}{@{}llrrr@{}}
    \toprule
    Suite & Task & Trial & Control steps & Terminal reward \\
    \midrule
    libero\_spatial & 4 & 4  & 600 & 0.0 \\
    libero\_spatial & 4 & 27 & 600 & 0.0 \\
    libero\_spatial & 5 & 2  & 600 & 0.0 \\
    libero\_object  & 7 & 10 & 600 & 0.0 \\
    libero\_object  & 7 & 13 & 600 & 0.0 \\
    libero\_goal    & 2 & 38 & 600 & 0.0 \\
    libero\_goal    & 2 & 47 & 600 & 0.0 \\
    libero\_goal    & 8 & 25 & 600 & 0.0 \\
    libero\_10      & 0 & 7  & 700 & 0.0 \\
    libero\_10      & 0 & 14 & 700 & 0.0 \\
    libero\_10      & 0 & 38 & 700 & 0.0 \\
    libero\_10      & 6 & 23 & 700 & 0.0 \\
    libero\_10      & 8 & 29 & 700 & 0.0 \\
    libero\_10      & 8 & 38 & 700 & 0.0 \\
    libero\_10      & 9 & 25 & 700 & 0.0 \\
    \bottomrule
  \end{tabular}
\end{table}

The failures span drawer retrieval, dairy placement, placing a bottle on a
cabinet, putting a pot on the stove and combined operations, but a task name is
not a diagnosed failure mechanism. Without the full video and action record of
each failing episode, the failures cannot be attributed to occlusion, memory,
grasping, planning or model hallucination. Any future diagnosis must register a
new run and new random conditions; new footage from a different random stream
cannot be presented as these 15 original failure trajectories.

\subsection{Reproducibility Notes}
\label{sec:owa-issues}

The first smoke attempt failed because the author implementation rejected the
\texttt{torch\_sdpa} attention enum, producing 0 episodes; after switching to
the legal \texttt{torch} value, a second attempt still produced 0 episodes
because \texttt{pandas} was missing during normaliser import. Both failures and
their logs were kept.

The root cause is that three packages eagerly re-export a training-only LeRobot
reader alongside the inference normaliser. Because the private disk lacked room
to install the full \texttt{pandas}/\texttt{pyarrow} training stack, only the
three initialisers were changed to lazy loading in a separate copy of the
author source; the remaining 546 files (including the model and the
normalisation maths) are SHA-identical, and a real CPU normaliser import and
min--max round trip passed. The tool is provided as
\btt{prepare\_inference\_imports.py}.
On 2026-10-06 the same 546 files and three initialisers were re-verified.

The teaching retry then succeeded (1/1, 74 steps) and the full run completed
once with exit code 0 and \btt{client\_max\_attempts=1};
no best-of-retry results were spliced into the main table. TCP readiness
probes during server start-up produced \texttt{InvalidMessage} handshake log
lines, which are not failed policy tasks --- log lines containing the word
``error'' are not automatically treated as NaN.

\subsection{Teaching Case: Images, Actions and Video}
\label{sec:owa-case}

The teaching case fixes Spatial task 0, initial state 0 and succeeded once in 74
control steps with 8 planning calls. All native 10-D end-effector, 7-D LIBERO
and 10-D state values are finite. It is not an extra 2001st episode of the main
table. The input images are the data the model genuinely received, not paper
screenshots or generated images.

The state is a 3-D position, a 6-D rotation and a 1-D gripper value
($3 + 6 + 1 = 10$); the internal 80-D representation is a unified model format,
not 80 robot degrees of freedom. The following record is the genuine first
control step of the teaching case. Note that the native action and the final
control differ in rotation and gripper due to adaptation, and that the final
7-D action is not the first seven entries of the native 10-D action.

\begin{verbatim}
{
  "state": [
    -0.21130245923995972, -0.010944196954369545, 1.1741831302642822,
    0.9984198808670044, -0.00016182575200218707, -0.056193094700574875,
    -0.0001631764171179384, -1.0, -1.944792347785551e-05, 0.975337028503418
  ],
  "native_eef10_action": [
    0.3204345405101776, 0.3149413764476776, 0.1583862155675888,
    0.99609375, 0.0078125, -0.0302734375, -0.0146484375,
    0.99609375, -0.0087890625, 1.0078125
  ],
  "libero_action7": [
    0.3204345405101776, 0.3149413764476776, 0.1583862155675888,
    -0.009145314805209637, 0.030346114188432693, 0.007979447022080421, -1.0
  ]
}
\end{verbatim}

Figure~\ref{fig:owa-case-sheet} shows the two genuine input views of the case,
and Figure~\ref{fig:owa-case-actions} plots the executed actions together with
the end-effector position.

\begin{figure}[H]
  \centering
  \includegraphics[width=0.62\linewidth]{2026.10.05-2026.10.11-openwam-libero/case-contact-sheet.png}
  \caption{Two genuine input views of the teaching case.}
  \label{fig:owa-case-sheet}
\end{figure}

\begin{figure}[H]
  \centering
  \includegraphics[width=0.9\linewidth]{2026.10.05-2026.10.11-openwam-libero/case-actions.png}
  \caption{Executed actions and end-effector position of the teaching case.}
  \label{fig:owa-case-actions}
\end{figure}

Figures~\ref{fig:owa-step0},~\ref{fig:owa-step36} and~\ref{fig:owa-step73}
show the two camera views at three points of the episode: before the first
action, at step 36 and at step 73. The saved frames precede their actions, and
success is decided by the post-action environment signal rather than by the
last image.

\begin{figure}[H]
  \centering
  \begin{subfigure}[t]{0.48\linewidth}
    \centering
    \includegraphics[width=\linewidth]{2026.10.05-2026.10.11-openwam-libero/step-0000-head.png}
    \caption{external view}
  \end{subfigure}\hfill
  \begin{subfigure}[t]{0.48\linewidth}
    \centering
    \includegraphics[width=\linewidth]{2026.10.05-2026.10.11-openwam-libero/step-0000-wrist.png}
    \caption{wrist view}
  \end{subfigure}
  \caption{Observation before the first action (step 0).}
  \label{fig:owa-step0}
\end{figure}

\begin{figure}[H]
  \centering
  \begin{subfigure}[t]{0.48\linewidth}
    \centering
    \includegraphics[width=\linewidth]{2026.10.05-2026.10.11-openwam-libero/step-0036-head.png}
    \caption{external view}
  \end{subfigure}\hfill
  \begin{subfigure}[t]{0.48\linewidth}
    \centering
    \includegraphics[width=\linewidth]{2026.10.05-2026.10.11-openwam-libero/step-0036-wrist.png}
    \caption{wrist view}
  \end{subfigure}
  \caption{Observation before the action at step 36.}
  \label{fig:owa-step36}
\end{figure}

\begin{figure}[H]
  \centering
  \begin{subfigure}[t]{0.48\linewidth}
    \centering
    \includegraphics[width=\linewidth]{2026.10.05-2026.10.11-openwam-libero/step-0073-head.png}
    \caption{external view}
  \end{subfigure}\hfill
  \begin{subfigure}[t]{0.48\linewidth}
    \centering
    \includegraphics[width=\linewidth]{2026.10.05-2026.10.11-openwam-libero/step-0073-wrist.png}
    \caption{wrist view}
  \end{subfigure}
  \caption{Observation before the last action (step 73).}
  \label{fig:owa-step73}
\end{figure}

The corresponding 74-step two-camera video is stored alongside the figures as
\texttt{case-rollout.mp4} and can be opened in a local player. It is the
genuine smoke case from before the full run, not an author demo or a generated
video, and not an original recording of the W2 episodes. The
\texttt{head\_camera} field is the simulator's external camera, not a
human-worn camera.

\subsection{Discussion: Reading OpenWAM Alongside \texorpdfstring{$\pi_{0.5}$}{pi-0.5}}
\label{sec:owa-discussion}

Table~\ref{tab:owa-vs-pi05} places the two reproductions side by side.

\begin{table}[H]
  \centering
  \small
  \caption{OpenWAM-$\alpha$ versus $\pi_{0.5}$, each under its own native
  protocol. The two columns answer separate reproduction questions.}
  \label{tab:owa-vs-pi05}
  \begin{tabular}{@{}lrr@{}}
    \toprule
    Item & $\pi_{0.5}$ (Exp.~B) & OpenWAM-$\alpha$ (Exp.~A) \\
    \midrule
    Mean success rate        & 96.75\,\% & 99.25\,\% \\
    Official mean            & 96.85\,\% & 99.30\,\% \\
    Offset from official     & $-0.10$\,pp & $-0.05$\,pp \\
    Environment seed         & 7   & 42 \\
    Stabilisation wait       & 10  & 30 \\
    Spatial/Object/Goal budget & 220/280/300 & 600/600/600 \\
    LIBERO-10 budget         & 520 & 700 \\
    Execution per plan       & predict 10, execute first 5 & native cache cadence, horizon 10 \\
    MuJoCo                   & 3.2.3 & 3.3.2 \\
    Per-query audit          & saved and independently verified & debug off, no equivalent record \\
    \bottomrule
  \end{tabular}
\end{table}

The two evaluations share the same four LIBERO suites, but the models, initial
random conditions, time budgets, control cadence and simulator versions were
not unified. The two columns therefore answer each author's own reproduction
question, and the difference $99.25\,\% - 96.75\,\% = 2.50$ percentage points
must not be presented as a fair architectural gain. A model-to-model comparison
would require a separate pre-registered protocol that fixes the environment,
seed, initial states, budget, action cadence and cost. The high score obtained
here is a baseline reproduction, not a new method.

\subsection{Limitations}
\label{sec:owa-limitations}

The claims of this experiment are bounded as follows. Evaluation is closed-loop
action success in simulation only; no real-robot deployment is attempted.
World-model video quality is not measured. Only one seed and one initial-state
set per task are used, so multi-seed stability is not established. The full run
keeps no per-query media, which limits post-hoc failure analysis. The
environment differs from the authors' recommendation in Python and Torch
versions. The official numbers are summary scores with no per-episode archive,
so exact agreement with the authors cannot be verified.

\subsection{Summary}
\label{sec:owa-summary}

On a single A100 with the authors' released weights and native protocol,
OpenWAM-$\alpha$ reproduces $1985/2000 = 99.25\,\%$ on the four LIBERO suites,
$0.05$ percentage points below the official 99.30\,\%. All 15 failures are
retained, the environment and provenance are documented, and a genuine teaching
case is provided with images, actions and a video. The result is a usable,
verifiable baseline of the authors' policy together with engineering and
reproducibility experience --- not a new contribution, and not automatically
transferable to human egocentric video.
